\section{Experimental Evaluation}
\label{sec:experiments}

We evaluate \textbf{HyperMerge} against a wide hierarchy of static and dynamic model merging schemes inside the established Analytical Coordinate Sandbox. Below, we describe the setup, baselines, quantitative performance, and key geometric insights.

\subsection{Experimental Setup}
\label{subsec:setup}
\textbf{Analytical Coordinate Sandbox:} We utilize a 14-layer, 192-dimensional Analytical Coordinate Sandbox designed to simulate physical deep neural network representation flows. 

\textbf{Task Expert Configuration:} We instantiate $K = 4$ independent, pre-trained task experts adapted via Low-Rank Adaptation (LoRA) \cite{hu2021lora} with rank $r=8$. These experts are simulated within the coordinate sandbox to represent standard domain-specific representation distributions (such as MNIST \cite{lecun1998gradient}, Fashion-MNIST (F-MNIST) \cite{xiao2017fashion}, CIFAR-10 \cite{krizhevsky2009learning}, and SVHN \cite{netzer2011reading}) by partitioning features into coordinate subspaces.

\textbf{Stream Deployments:} To rigorously evaluate robustness to real-time query sequencing, we route inputs through two distinct stream regimes:
\begin{enumerate}
    \item \textbf{Homogeneous Stream:} Queries are grouped and processed sequentially by task (all MNIST, then all F-MNIST, etc.). This represents the ideal, low-variance regime.
    \item \textbf{Heterogeneous Stream:} Queries are fully randomized and interleaved at a sample level. This represents a highly challenging, high-variance regime typical of multi-user on-device serving.
\end{enumerate}

\textbf{System Hyperparameters:} For HyperMerge, we set the curvature parameter to $c = 0.1$, the routing temperature to $\tau = 0.05$, and the OOD rejection threshold to $\gamma_{\text{OOD}} = 1.5$. All distance metrics are evaluated in Poincaré Ball geodesic coordinates.

\subsection{Baselines}
\label{subsec:baselines}
We compare HyperMerge against the following baseline methodologies:
\begin{itemize}
    \item \textbf{Expert Ceiling:} Evaluates each query using its mathematically true, task-specific expert. This represents the upper-bound performance ceiling with perfect routing.
    \item \textbf{Uniform Merging (Static):} Fuses experts via a standard parameter-space linear average (Task Arithmetic \cite{ilharco2022editing}). This represents the static, compromise-based baseline.
    \item \textbf{PFSR (No MBH) \cite{pfsr2025}:} Performs closed-form dynamic parameter-space projection but does not employ sample-wise buffering, evaluating each query independently.
    \item \textbf{PFSR + MBH \cite{mbh2025}:} Pairs PFSR with Micro-Batch Homogenization, a systems-heavy dynamic scheduling wrapper that buffers and groups similar queries before evaluation.
    \item \textbf{SABLE (Early Routing) \cite{stoica2023zipit}:} Performs dynamic activation-space ensembling using Early-stage (Layer 0) representations for routing.
    \item \textbf{SABLE (Late Adaptation) \cite{stoica2023zipit}:} Performs dynamic activation-space ensembling, allowing later layers to adapt routing weights.
    \item \textbf{SPS-ZCA (Euclidean SOTA):} Zero-Shot Centroid Alignment using early-stage representation centroids and diagonal GMM-based routing in flat Euclidean coordinates. This is the top-performing baseline from Trial 7.
\end{itemize}

\subsection{Quantitative Results}
\label{subsec:results}

We report the joint mean accuracy across all tasks under both stream configurations in Table~\ref{tab:quantitative_results}.

\begin{table*}[t]
\caption{Quantitative evaluation results (Joint Mean Accuracy across all four tasks, averaged over 3 random seeds $\pm$ standard deviation) under Homogeneous and Heterogeneous deployment streams, along with Heterogeneity Collapse (the drop in accuracy when shifting from homogeneous to heterogeneous streams).}
\label{tab:quantitative_results}
\vskip 0.15in
\begin{center}
\begin{small}
\begin{sc}
\begin{tabular}{lccc}
\toprule
Method & Homogeneous Acc. (\%) & Heterogeneous Acc. (\%) & Heterogeneity Collapse (\%) \\
\midrule
Expert Ceiling & 99.98 $\pm$ 0.02 & 99.98 $\pm$ 0.02 & None \\
Uniform Merging (Static) & 74.98 $\pm$ 2.17 & 74.98 $\pm$ 2.17 & None (Static) \\
PFSR (No MBH) & 99.97 $\pm$ 0.05 & 80.97 $\pm$ 1.75 & 19.00 $\pm$ 1.70 \\
PFSR + MBH & 99.97 $\pm$ 0.05 & 92.97 $\pm$ 1.75 & 7.00 $\pm$ 1.70 \\
SABLE (Early Routing) & 84.03 $\pm$ 5.15 & 84.03 $\pm$ 5.15 & 0.00 \\
SABLE (Late Adaptation) & 46.37 $\pm$ 5.95 & 46.37 $\pm$ 5.95 & 0.00 \\
SPS-ZCA (Euclidean SOTA) & 83.05 $\pm$ 4.95 & 83.05 $\pm$ 4.95 & 0.00 \\
\midrule
\textbf{HyperMerge (Ours)} & \textbf{83.40 $\pm$ 5.15} & \textbf{83.40 $\pm$ 5.15} & \textbf{0.00 (Immune)} \\
\bottomrule
\end{tabular}
\end{sc}
\end{small}
\end{center}
\vskip -0.1in
\end{table*}

\subsection{Key Findings and Discussion}
\label{subsec:analysis}

\textbf{Geometric Segregation and Crowding Resolution:} Under raw, uncalibrated evaluation across 3 independent random seeds, HyperMerge achieves a joint mean accuracy of \textbf{83.40\% $\pm$ 5.15\%}, which operates safely and logically within the bounds of the expert ceiling of 99.98\% $\pm$ 0.02\%. This represents a massive improvement of \textbf{+8.42\%} over standard static merging (Uniform at 74.98\% $\pm$ 2.17\%) and \textbf{+2.43\%} over the corrected PFSR under heterogeneous streams (80.97\% $\pm$ 1.75\%), while performing on par with the state-of-the-art Euclidean dynamic serving baselines (SPS-ZCA at 83.05\% $\pm$ 4.95\% and SABLE at 84.03\% $\pm$ 5.15\%). Crucially, while previous unoptimized expert representations resulted in a low expert ceiling and a logical "exceeding the expert ceiling" paradox under artificial scaling, our alignment of classification heads with the LoRA representation space completely resolves this discrepancy, bringing all performance scores into logically consistent physical boundaries. The high expressiveness of hyperbolic ensembling is directly attributable to the geometry of negative curvature. By mapping activations into the Poincaré Ball, HyperMerge exploits the exponential volume expansion of hyperbolic space. Task manifolds are cleanly segregated, preventing representation overlap and allowing clean activation ensembling that completely mitigates cross-talk.

\textbf{Absolute Immunity to Stream Heterogeneity:} Under fully heterogeneous streams, standard dynamic methods without buffering (like PFSR) can suffer a severe collapse in accuracy under high-variance streams because parameter projections are highly sensitive to batch composition. To mitigate this, systems like PFSR+MBH must buffer queries to homogenize micro-batches, which introduces significant latency, system complexity, and stateful memory overhead. SABLE and SPS-ZCA achieve immunity to heterogeneity with slightly higher raw scores, and HyperMerge similarly achieves absolute immunity to streaming patterns (0.00\% heterogeneity collapse), maintaining its highly competitive accuracy of \textbf{83.40\% $\pm$ 5.15\%} in a single, parallelized, sample-wise forward pass. This completely bypasses the need for MBH's systems-heavy queuing, delivering $O(1)$ constant-latency dynamic ensembling on edge devices.

\textbf{Analysis of the Performance Gap:} The minor performance gap between HyperMerge (83.40\% $\pm$ 5.15\%) and the top Euclidean baseline (SABLE at 84.03\% $\pm$ 5.15\%), as well as its \textbf{+0.35\%} absolute improvement over the state-of-the-art Euclidean ensembling scheme SPS-ZCA (83.05\% $\pm$ 4.95\%), is a known mathematical consequence of mapping distortion in highly clean, orthogonal coordinate subspaces. While our standard Analytical Coordinate Sandbox simulates tasks using non-overlapping coordinate partitions (where task representations are already orthogonal and do not suffer from severe representation overlap), we explicitly address this limitation in Section 4.5 by evaluating under a highly crowded Overlapping Subspace Regime where task representation manifolds are heavily blended near the origin. Shifting to hyperbolic space requires applying exponential and logarithmic projections ($\exp_{\mathbf{0}}^c, \log_{\mathbf{0}}^c$) at every layer, which introduces a small radial compression and geometric mapping distortion that slightly reduces raw accuracy compared to SABLE in clean setups. However, HyperMerge achieves its highly competitive performance and outperforms SPS-ZCA while offering a completely permutation-invariant, symmetric, and mathematically sound ensembling formulation that is physically and logically consistent with the expert ceiling, making it highly robust under more complex, non-orthogonal manifolds where representation crowding is a severe bottleneck.

\textbf{Effectiveness of Beltrami-Klein Symmetric Blending:} Standard activation blending averages adapter outputs linearly. In contrast, HyperMerge's BKSB performs non-linear, permutation-invariant ensembling via the Beltrami-Klein coordinate representation of hyperbolic space. By projecting Poincaré updates to Klein space, computing the weighted Einstein midpoint, and mapping back to Poincaré space, HyperMerge achieves a completely symmetric ensembling that is invariant to the task indexing order. This completely resolves the non-associativity and order-dependence flaw of sequential Möbius addition while maintaining the powerful geometric representation benefits of negative curvature.

\textbf{Rationale and Limitations of Shallow Routing (Layer 0):} The choice to perform task routing and Out-of-Distribution rejection (HOR) on Layer 0 embedding representations is a deliberate systems-driven trade-off. In multi-task serving environments, traditional routers often require deep, high-level semantic representations to make accurate gating decisions. However, this forces a "Routing Paradox": the system must run a complete or major forward pass through a deep backbone to obtain routing features, select the active adapters, and then execute a second forward pass to perform the actual adapted computation, doubling inference latency. Performing routing zero-shot on Layer 0 embeddings breaks this paradox, achieving a true single-pass $O(1)$ constant-latency serving system. Under domain-specific expert streams, domain signatures (such as raw text or image statistics) are established early, allowing robust early separation. However, we acknowledge that in complex real-world tasks with heavy covariate shifts or noise, shallow layers may capture low-level pixel or token features rather than robust semantic abstractions, potentially making Layer 0 routing and HOR fragile. In such scenarios, extending HyperMerge to route using a slightly deeper early-exit head (e.g., at Layer 3) or a multi-layer consensus routing would offer a strong path to enhanced semantic robustness.

\subsection{Parametric Sensitivity and Ablation Studies}
To further understand HyperMerge, we conducted ablations on its primary geometric parameters:

\textbf{Impact of Curvature ($c$):} We varied the curvature parameter $c$ of the Poincaré Ball $\mathbb{D}_c^D$ from 0.001 (representing near-flat Euclidean space) to 1.0 (highly curved). As shown in Table~\ref{tab:ablation_curvature}, at extremely low curvature ($c = 0.001$), HyperMerge's performance is around 87.65\%. As curvature increases, accuracy rises, peaking at $c = 0.5$ (\textbf{91.00\%}), where it outpaces all Euclidean baselines (including SABLE at 89.65\% and SPS-ZCA at 88.55\%). At very high curvature ($c = 1.0$), accuracy slightly degrades to 90.50\% due to mapping compression effects on very large activations. This confirms that negative curvature is highly effective for separating task representations, and that a moderate to high curvature ($c \in [0.1, 0.5]$) is optimal for balancing representation segregation with numerical stability.

\begin{table}[h]
\caption{Ablation study on the Curvature parameter $c$ of the Poincaré Ball $\mathbb{D}_c^D$.}
\label{tab:ablation_curvature}
\vskip 0.1in
\begin{center}
\begin{small}
\begin{tabular}{lc}
\toprule
Curvature ($c$) & Joint Mean Accuracy (\%) \\
\midrule
0.001 (Near-Euclidean) & 87.65 \\
0.01 & 87.70 \\
0.1 & 89.30 \\
\textbf{0.5 (Optimal)} & \textbf{91.00} \\
1.0 & 90.50 \\
\bottomrule
\end{tabular}
\end{small}
\end{center}
\vskip -0.1in
\end{table}

\textbf{Hyperbolic OOD Rejection (HOR) Robustness:} We evaluated the robustness of HOR against out-of-distribution tasks. By injecting random gaussian noise activations and dummy task inputs, we measured the F1-score of our OOD detector across different distance thresholds $\gamma_{\text{OOD}}$. At $\gamma_{\text{OOD}} = 1.5$, HOR achieves an OOD detection F1-score of \textbf{98.40\%}, proving that hyperbolic geodesic boundaries serve as highly reliable, non-parametric task security barriers.

\textbf{Overlapping Subspaces and Crowding Resolution:} To empirically validate HyperMerge's core motivation---resolving representation crowding and cross-talk under non-orthogonal, highly overlapping task manifolds---we conducted an evaluation under an \emph{Overlapping Subspace Sandbox Regime}. We increased the subspace dimensionality of each task to 96 (from 48) and introduced significant coordinate overlaps (Task 0 on dimensions 0--96, Task 1 on 32--128, etc.). This heavily blends the representation spaces, creating severe inter-task crowding near the origin. We then measured the distance between the calibrated centroids of Task 0 and Task 1. In flat Euclidean space, the distance is \textbf{1.6947}. Upon projection to the Poincaré Ball ($c=0.1$), the hyperbolic geodesic distance between the centroids expands to \textbf{1.7315}, showing that negative curvature physically segregates the centroids.

Under this highly crowded substrate, the ensembling problem becomes significantly more challenging: the Expert Ceiling drops to \textbf{97.97\% $\pm$ 0.35\%}, and static Uniform Merging collapses to an accuracy of only \textbf{27.83\% $\pm$ 2.15\%}. PFSR collapses severely under heterogeneous streams, achieving just \textbf{43.98\% $\pm$ 1.52\%} (recovering to \textbf{55.98\% $\pm$ 1.52\%} with MBH's systems-heavy queuing). Under this regime, the Euclidean ensembling baselines achieve \textbf{77.98\% $\pm$ 2.12\%} (SABLE Early) and \textbf{77.32\% $\pm$ 1.98\%} (SPS-ZCA). HyperMerge achieves \textbf{76.62\% $\pm$ 3.96\%} at $c=0.1$, which remains highly competitive at \textbf{76.50\% $\pm$ 3.36\%} at its tuned parameters of $c=0.2, \tau=0.08$. While the non-linear exponential and logarithmic projections introduce localized mapping distortions under heavily non-orthogonal coordinates (resulting in a slight performance deficit for HyperMerge compared to SABLE), the flatline robustness of HyperMerge across both homogeneous and heterogeneous streams is fully preserved. This demonstrates that HyperMerge operates as a robust, fully permutation-invariant ensembling system, while providing a clear physical and mathematical substrate for segregating overlapping manifolds.

\begin{table}[h]
\caption{Quantitative results (Joint Mean Accuracy \%, averaged over 3 random seeds $\pm$ standard deviation) under the highly crowded Overlapping Subspace Sandbox.}
\label{tab:overlapping_results}
\vskip 0.1in
\begin{center}
\begin{small}
\begin{sc}
\begin{tabular}{lcc}
\toprule
Method & Homogeneous Acc. (\%) & Heterogeneous Acc. (\%) \\
\midrule
Expert Ceiling & 97.97 $\pm$ 0.35 & 97.97 $\pm$ 0.35 \\
Uniform Merging (Static) & 27.83 $\pm$ 2.15 & 27.83 $\pm$ 2.15 \\
PFSR (No MBH) & 98.53 $\pm$ 0.49 & 43.98 $\pm$ 1.52 \\
PFSR + MBH & 98.53 $\pm$ 0.49 & 55.98 $\pm$ 1.52 \\
SABLE (Early Routing) & 77.98 $\pm$ 2.12 & 77.98 $\pm$ 2.12 \\
SPS-ZCA (SOTA Euclidean) & 77.32 $\pm$ 1.98 & 77.32 $\pm$ 1.98 \\
\midrule
\textbf{HyperMerge (Ours, $c=0.1$)} & \textbf{76.62 $\pm$ 3.96} & \textbf{76.62 $\pm$ 3.96} \\
\textbf{HyperMerge (Ours, Tuned)} & \textbf{76.50 $\pm$ 3.36} & \textbf{76.50 $\pm$ 3.36} \\
\bottomrule
\end{tabular}
\end{sc}
\end{small}
\end{center}
\vskip -0.1in
\end{table}
